\begin{table}[!htbp]
\centering
\small
\setlength{\tabcolsep}{3.5pt}
\renewcommand{\arraystretch}{1.04}
\resizebox{\textwidth}{!}{%
\begin{tabular}{@{}>{\raggedright\arraybackslash}p{0.37\textwidth}rr@{}}
\toprule
Simulation setting & AL & exAL \\
\midrule
Asymmetric-Laplace tail & -0.0355 [-0.0649, -0.0060] & -0.0586 [-0.0954, -0.0242] \\
Gaussian-mixture innovations & 0.0090 [-0.0168, 0.0348] & 0.0136 [-0.0164, 0.0433] \\
Laplace innovations & 0.1529 [0.0823, 0.2397] & 0.1430 [0.0508, 0.2688] \\
Nonlinear reservoir dynamics & 0.0120 [-0.0098, 0.0392] & 0.0034 [-0.0212, 0.0280] \\
Gaussian innovations & -0.0143 [-0.0368, 0.0062] & 0.0010 [-0.0237, 0.0254] \\
Persistent heavy tails & 0.0695 [-0.2180, 0.3848] & -0.0233 [-0.2876, 0.2513] \\
Regime shift & 0.1024 [-0.2039, 0.4033] & 0.0087 [-0.2867, 0.2706] \\
Student-\(t\) location--scale & 0.0010 [-0.0020, 0.0040] & 0.0013 [-0.0035, 0.0060] \\
\bottomrule
\end{tabular}
}%
\caption{Joint-minus-independent differences in DGP-integrated \(\aCRPS\), with means and equal-tailed 95\% intervals. Positive differences favor independent estimation. The deterministic index pairing defines descriptive contrasts between separately fitted models; the two intervals for the asymmetric-tail setting favor joint estimation, the two Laplace intervals favor independent estimation, and twelve include zero.}
\label{tab:joint-qdesn-pure-desn-v1-contrasts}
\end{table}
